\chapter{BMI and Waist Circumference}

\section*{Definition and Overview}
Body mass index (BMI) and waist circumference are simple measurements used to assess whether a person's body weight and fat distribution place them at higher risk of health problems. BMI is calculated from a person's weight and height, while waist circumference measures the amount of fat around the abdomen. Used together, they give a fuller picture of health risk than either measurement alone, because carrying extra weight around the middle is a particularly strong marker of risk for conditions such as type 2 diabetes, high blood pressure, and cardiovascular disease.

\section*{Epidemiology}
Overweight and obesity are common and increasing worldwide. Roughly one in four adults in many developed countries has obesity, and about half the adult population is above the healthy weight range. Abdominal obesity is even more widespread, and the rise is seen in children and adolescents as well as adults. Excess weight is a major contributor to diabetes, heart disease, stroke, some cancers, and joint problems, making BMI and waist measurement important tools for identifying people who would benefit from weight-related health support.

\section*{Aetiology and Pathophysiology}
BMI and waist circumference work because body fat affects health:

\begin{itemize}
    \item \textbf{BMI:} weight in kilograms divided by height in metres squared -- a healthy range is roughly 18.5--24.9, with 25--29.9 considered overweight and 30 or above considered obese
    \item \textbf{Waist circumference:} measured around the abdomen at the level of the navel; a high waist measurement (roughly above 94 cm for men and 80 cm for women, with ethnic variation) indicates excess abdominal fat
    \item \textbf{Abdominal fat:} fat stored around the organs is metabolically active, releasing substances that promote insulin resistance, inflammation, and high blood pressure
    \item \textbf{Limits of BMI:} it does not distinguish muscle from fat, so a very muscular person may have a high BMI without excess fat
    \item \textbf{Energy balance:} weight gain results when energy taken in exceeds energy used over a long period
\end{itemize}

\section*{Clinical Features}
The measurements themselves cause no symptoms. What they point to are the consequences of excess weight:

\begin{itemize}
    \item A weight above the healthy range or a large waist measurement
    \item Often no symptoms at all in the early stages
    \item Possible tiredness, breathlessness on exertion, snoring, or joint aches as weight increases
    \item Associated conditions may be present without being noticed, such as raised blood pressure, abnormal cholesterol, or high blood sugar
    \item Changes in appearance and self-esteem can affect well-being
\end{itemize}

\section*{Differential Diagnosis}
A high BMI is not the same as being unhealthy, and other factors should be considered:

\begin{itemize}
    \item Muscle mass can raise BMI in athletic people without excess fat
    \item Fluid retention, pregnancy, or body build affect weight and measurements
    \item Underlying conditions such as an underactive thyroid or medication side effects may contribute to weight gain
    \item Weight can be normal despite an unhealthy diet or physical inactivity
    \item A person with a normal BMI can still have abdominal obesity or metabolic risk factors
\end{itemize}

\section*{When to Seek Medical Care}
See a doctor if:

\begin{itemize}
    \item You are unsure what your BMI or waist circumference means for your health
    \item You have risk factors such as diabetes, high blood pressure, heart disease, or a strong family history
    \item You have symptoms such as breathlessness, snoring, daytime sleepiness, or joint pain
    \item Your weight has changed rapidly without a clear reason
    \item You have struggled with weight and want a structured, supported approach
\end{itemize}

\textbf{Call 000 (or local emergency services) for:}
\begin{itemize}
    \item Chest pain, difficulty breathing, or sudden weakness or numbness
    \item Any sudden, severe symptom that worries you
\end{itemize}

\section*{Diagnosis and Investigations}
\begin{itemize}
    \item \textbf{Height and weight:} measured with proper equipment to calculate BMI
    \item \textbf{Waist circumference:} measured with a tape around the abdomen at the navel level
    \item \textbf{Blood pressure:} raised blood pressure is common with excess weight
    \item \textbf{Blood tests:} may include blood sugar, cholesterol, liver function, and thyroid function
    \item \textbf{Body composition:} more detailed methods, such as bioimpedance or DEXA scans, are sometimes used to measure fat versus muscle
    \item \textbf{History:} diet, physical activity, family history, medicines, and sleep patterns
\end{itemize}

\section*{Management}
Weight management is a long-term effort best supported by a team approach:

\begin{itemize}
    \item \textbf{Realistic goals:} a modest weight loss of 5--10\% produces meaningful health benefits, even if the target weight is not reached
    \item \textbf{Diet:} a balanced eating pattern, reducing portion sizes, sugary drinks, and highly processed foods, and increasing vegetables, fruit, and whole grains
    \item \textbf{Physical activity:} building up to regular aerobic activity plus strengthening exercises
    \item \textbf{Behavioural support:} structured programs, goal setting, and self-monitoring improve results
    \item \textbf{Medication:} weight-loss medicines are available and are used under medical supervision as part of a broader program
    \item \textbf{Surgery:} bariatric surgery is an option for people with severe obesity, after careful assessment
    \item \textbf{Treating complications:} managing diabetes, blood pressure, and sleep apnoea alongside weight loss
\end{itemize}

\section*{Complications}
Excess weight and abdominal fat increase the risk of:

\begin{itemize}
    \item Type 2 diabetes and metabolic syndrome
    \item Heart disease, high blood pressure, and stroke
    \item Fatty liver disease and gallbladder problems
    \item Certain cancers, including bowel, breast, and kidney cancer
    \item Obstructive sleep apnoea and breathing problems
    \item Osteoarthritis, back pain, and reduced mobility
    \item Lower self-esteem, depression, and reduced quality of life
\end{itemize}

\section*{Prognosis and Outcomes}
Weight loss of even 5--10\% often improves blood pressure, blood sugar, cholesterol, and sleep quality, and these improvements appear regardless of whether a person reaches a healthy BMI. Weight regain is common, and sustained success usually depends on habits that can be maintained long term rather than short-term dieting. People who are supported by health professionals and who set realistic goals are more likely to keep the weight off. The overall risk of diabetes and heart disease falls substantially with lasting weight reduction.

\section*{Prevention and Self-Care}
\begin{itemize}
    \item Aim for a healthy weight range for your height, with a waist circumference below the high-risk threshold
    \item Eat regular meals with plenty of vegetables and fruit, and limit sugary drinks and snacks
    \item Be active most days, combining aerobic activity with strength exercises
    \item Watch portion sizes and avoid eating late at night out of boredom or stress
    \item Get enough good-quality sleep, since poor sleep is linked to weight gain
    \item Weigh yourself regularly, but do not let daily fluctuations cause distress
    \item Seek support early -- from a GP, dietitian, or structured program -- rather than struggling alone
\end{itemize}

\section*{Sources and Further Reading}
The following sources provide reliable, up-to-date information on BMI, waist circumference, and weight management:

\begin{itemize}
    \item World Health Organization. \textit{Obesity and overweight fact sheets.} https://www.who.int/
    \item NHS. \textit{NHS guidance on BMI, weight, and healthy living.} https://www.nhs.uk/conditions/
    \item Mayo Clinic. \textit{Information on body mass index and weight management.} https://www.mayoclinic.org/
    \item Centers for Disease Control and Prevention. \textit{About adult BMI and healthy weight.} https://www.cdc.gov/
    \item MedlinePlus. \textit{Body mass index patient information.} https://medlineplus.gov/
    \item NICE. \textit{Obesity identification, assessment, and management guidelines.} https://www.nice.org.uk/
\end{itemize}
